\section{Affine motion and exceptional physical time}
\label{sec:motion}

The motion obstruction concerns actual work-head trajectories, including
setup and cleanup.  We derive a quantitative consequence of the first paper's
limiting inequality \cite{FirstPaper2026}, then transfer it to native
executions.  This provides the lower-bound criterion used to exclude the
concrete palettes and recursive schedules in Section~\ref{sec:motion-applications}.
The first paper remarks that a common triangle wave fails its inequality;
it does not contain an affine-motion theorem.  The results below use only
the $S=\varnothing$, $0<\theta<1/2$ instances of that inequality.

\subsection{The crossing convention and the affine-time bound}

Write $\lambda$ for Lebesgue measure.  For a path $q:[0,H]\to\mathbb R$ and a
height $p$, let $\Lambda_q(p)=\{u\in[0,H]:q(u)=p\}$.  Two visits $u<v$ are
consecutive if no member of this \emph{full} level set lies strictly between
them.  For an excision set $S$ and a label $\beta:\mathbb R\to\mathbb N$, let
$\kappa_q(p;S,\beta)$ count consecutive pairs with both endpoints outside $S$
and with different labels.  As in the first paper, this count is defined to
be zero when $\Lambda_q(p)$ is infinite.  In particular, removing $S$ does not
create new consecutive pairs.

A label is admissible off $S$ if there is a finite set $Z$ such that it is
constant on every closed subinterval of $[0,H]\setminus(S\cup Z)$.  The values
at the finitely many breakpoints are unrestricted.  All crossing integrals
below are nonnegative Lebesgue integrals, with values in $[0,\infty]$.

\begin{definition}[The restricted star condition]
\label{def:motion-star}
For $k$ paths $q_t:[0,H]\to\mathbb R$, write $\mathcal S(q,H)$ if, for every
$0<\theta<1/2$, every finite union $S$ of closed intervals, and every admissible
label $\beta$ whose classes off $S$ have measure at most $\theta$,
\begin{equation}
\label{eq:motion-star}
  1\le \sum_{t=1}^k\int_{\mathbb R}\kappa_{q_t}(p;S,\beta)\,dp
       +2\theta+2k\lambda(S\cap[0,H]).
\end{equation}
A set $R_t\subseteq[0,H]$, open in $\mathbb R$, is an affine regular region
for $q_t$ if every point of $R_t$ has an open neighborhood contained in $R_t$
on which $q_t$ is affine.  The affine coefficients may depend on the
neighborhood and the head.
\end{definition}
\formal{FiniteAffineTimeBound}{TransposeFollowup.AffineMotion.SatisfiesStar}

The original $M^*$ conclusion implies $\mathcal S$; this restriction suffices
for the obstruction below.  The padding theorem retains the full original
condition.  Second countability gives each regular region a countable open
affine cover, so countability is not an additional assumption.  The cover
sets need not be intervals or form a partition.

\begin{theorem}[Quantitative affine-time obstruction]
\label{thm:motion-affine-time}
Let $H\ge0$, let each $q_t$ be $1$-Lipschitz on $[0,H]$, and assume
$\mathcal S(q,H)$.  For any affine regular regions $R_1,\ldots,R_k$,
\begin{equation}
\label{eq:motion-half}
  \sum_{t=1}^k\lambda([0,H]\setminus R_t)\ge\frac12.
\end{equation}
The same conclusion holds when each $R_t$ is open and $q_t$ is affine on each
connected component of $R_t$.  If $k>0$ and a common region $R$ works for every
head, then $\lambda([0,H]\setminus R)\ge1/(2k)$.
\end{theorem}
\formal{OpenAffineTimeBound}{TransposeFollowup.AffineMotion.star_half_le_locally_affine_omission,TransposeFollowup.AffineMotion.star_half_le_affine_component_omission,TransposeFollowup.AffineMotion.star_common_affine_omission}

The property \eqref{eq:motion-half}, quantified over all affine regular
regions, is named \texttt{AffineTimeRequired} in Lean.
\formal{NativeAffineMotion}{TransposeFollowup.AffineMotion.AffineTimeRequired}

\begin{proofoutline}
The labeling argument proves the following intermediate statement, which
does not require $\mathcal S(q,H)$.
For countable open affine covers $R_t=\bigcup_j U_{t,j}$, and any positive
$\theta,\varepsilon$, there is \emph{one} admissible measurable label $\beta$
with class measures at most $\theta$ such that
\begin{equation}
\label{eq:motion-small-cut}
 \sum_t\int\kappa_{q_t}(p;\varnothing,\beta)\,dp
 <2\sum_t\lambda([0,H]\setminus R_t)+\varepsilon.
\end{equation}
For finitely many nonconstant affine branches, inverse visits at a common
height are affine functions of that height.  Bound retained crossing pairs
by all ordered pairs of inverse branches and push their height measures to
pairs of times.  The resulting measures have explicit marginal bounds.
The affine group is solvable and hence amenable: its F\o lner property
provides finite sets of transformations with arbitrarily small relative
boundary under any prescribed finite family \cite{Folner1955}.
Our formal proof constructs these sets directly, first choosing slopes and
then translations.  A fine grid and a common priority labeling of those
sets make both label classes and the measure of disagreeing pairs small.
Grid collisions and finite boundaries are charged separately.  The finitely
many constant-branch heights have measure zero; their time intervals are
retained.

For a finite initial subfamily of the open cover, the complement in $[0,H]$
is compact.  Deleting its visits
can affect at most two consecutive pairs per deleted visit.  A measurable
visit-count majorant for a $1$-Lipschitz path bounds the integrated charge by
twice the complement's length.  This uses the original full-level adjacency
relation, without asserting that a clipped path has the same crossings.
Finite-prefix complements decrease to the actual omitted set, so continuity
of finite measure proves \eqref{eq:motion-small-cut} for countably many
branches.  Substituting the common label into \eqref{eq:motion-star} with
$S=\varnothing$, then sending $\theta,\varepsilon$ to zero, yields
\eqref{eq:motion-half}.  A countable cover of any locally affine open region
finishes the general case.
\end{proofoutline}
\formal{CountableAffinePathLabels}{TransposeFollowup.AffineMotion.exists_countable_affine_path_cut_bound}
\formal{FiniteAffineLabels}{TransposeFollowup.AffineMotion.exists_finite_affine_small_cut}
\formal{OmittedVisitBound}{TransposeFollowup.AffineMotion.lintegral_levelCross_delete_le}
\formal{AffineFiniteBoundary}{TransposeFollowup.FiniteBoundary.exists_affine_small_boundary}

Theorem~\ref{thm:motion-affine-time} quantifies the time a candidate must
leave outside its affine regions, even with countably many regions and
slopes of either sign, including zero.  The bound sums head-specific omitted
measures; it does not require a common exceptional set for all heads.
The F\o lner step connects the argument to the classical obstruction to
expansion from amenability; see Lubotzky--Weiss \cite{LubotzkyWeiss1993}.
The crossing-measure estimates and native time normalization are the
additional steps needed here.

\subsection{Native normalization and correct-input quantifiers}

For the native results, let $d_i\to\infty$. After discarding a finite prefix,
assume $d_i\ge1$. Put
$b_i=\max(1,\lceil\log_2 d_i\rceil)$, and $P_i=d_i^2b_i$.  Let $M_i$ be
native finite multitape machines with a fixed number $k$ of work tapes and a
fixed cell charge $G=\operatorname{cellBits}(M_i)\ge1$.  If $Q_i$ is the number
of finite-control states, assume
\begin{equation}
\label{eq:motion-control}
  \frac{\lceil\log_2 Q_i\rceil}{P_i}\longrightarrow0.
\end{equation}
Let $W_i$ be correct-input sets and $m_i\in\mathbb N$ satisfy, eventually,
\begin{equation}
\label{eq:motion-entropy}
  2^{P_i}\le 2^{m_i}|W_i|,\qquad m_i/P_i\longrightarrow0.
\end{equation}
Every member of $W_i$ must be the raw row-major binary encoding of a
$d_i\times d_i$ matrix, and the machine must actually halt with the raw
encoding of its transpose.  When we assume a fast execution, this halting
must occur by $cP_i$ ticks for one constant $c$.  All these input and time
conditions may hold only eventually in $i$.  There is no condition on inputs
outside $W_i$, and the machines themselves may vary with $i$.

Let $h_{i,t}^{w}(m)\in\mathbb Z$ be the position of work head $t$ after $m$
actual ticks on input $w$.  The native trajectory is frozen after halting.
Physical time and space are measured in units of $P_i$:
\begin{equation}
\label{eq:motion-physical}
  f_{i,t}^{w}(s)=\frac{\widetilde h_{i,t}^{w}(P_i s)}{P_i},
  \qquad
  \widehat f_{i,t}^{w}(s)=\frac{h_{i,t}^{w}(\lfloor P_i s\rfloor)}{P_i},
\end{equation}
where $\widetilde h$ is any $1$-Lipschitz extension agreeing at integer ticks.
Such an extension exists, and the two expressions differ by at most $1/P_i$
for $s\ge0$.

The block normalization in \cite{FirstPaper2026} is different at finite size.
Put $r_i=\lceil d_i b_i/G\rceil$.  Its actual block trajectory is
\[
  B_{i,t}^{w}(u)=\frac1{d_i}
     \left\lfloor\frac{h_{i,t}^{w}(r_i\lfloor d_i u\rfloor)}{r_i}\right\rfloor,
\]
where the outer floor is the integer block index, also for negative head
positions.  The proved comparison is
\begin{equation}
\label{eq:motion-normalization-error}
 \left|f_{i,t}^{w}(u/G)-B_{i,t}^{w}(u)/G\right|
 \le \frac{2(u+1)}{d_i}\quad(u\ge0).
\end{equation}
Thus a block limit $q_t$ corresponds to the physical limit
$f_t(s)=q_t(Gs)/G$, although the finite trajectories differ.

The next theorem assembles the cited native-limit and scaling declarations;
it is not a single exported Lean statement.

\begin{theorem}[Physical units and native limiting paths]
\label{thm:motion-native-limits}
Under \eqref{eq:motion-control}--\eqref{eq:motion-entropy} and fast correct
halting, there are representatives $w_i\in W_i$ eventually and a subsequence
whose block trajectories converge uniformly on $[0,cG+1]$.  Every uniform
subsequential limit of these representatives is $1$-Lipschitz and satisfies
\eqref{eq:motion-half}.  For the corresponding physical limits on
$[0,c+1/G]$, any affine regular regions satisfy
\begin{equation}
\label{eq:motion-physical-budget}
  \sum_{t=1}^k\lambda([0,c+1/G]\setminus R_t)\ge\frac1{2G}.
\end{equation}
When $k>0$, a common open region on each of whose connected components every
head is affine omits at least $1/(2kG)$ physical time.
The assertion includes a fixed uniform machine as a special case.
\end{theorem}
\formal{NativeAffineMotion}{TransposeFollowup.AffineMotion.popular_affine_time_nonuniform,TransposeFollowup.AffineMotion.popular_affine_time_robust,TransposeFollowup.AffineMotion.AffineTimeRequired.physical_omission}
\formal{AffineTimeScaling}{TransposeFollowup.AffineMotion.star_physical_common_affine_omission}
\formal{NativePhysicalLimits}{TransposeFollowup.RegularMotion.native_sampled_physical_limit}
\formal{CanonicalTimeNormalization}{TransposeFollowup.RegularMotion.canonical_physical_error}
\upformal{Transpose.Tape.MStarRobust}{Transpose.Tape.Impl.theoremMStar_level_nonuniform}

\begin{proofoutline}
The popular-trace theorem of \cite{FirstPaper2026} supplies actual trace
classes $F_i\subseteq W_i$ with representatives $w_i$, common halting times,
and integer deficits $D_i$ satisfying
$D_i\le m_i+C_0d_i b_i$ eventually, for a fixed trace-count constant $C_0$.
Consequently $D_i/P_i\to0$.  Its native compactness and $M^*$ theorem applies
to these classes; Theorem~\ref{thm:motion-affine-time} then applies to every
eligible limit.  The estimate \eqref{eq:motion-normalization-error} transfers
uniform convergence to physical paths, including all phases of execution
and the frozen halting tail.  Finally the change of variables $u=Gs$ scales
omitted measure by $G$.  For the common-region clause, retain the original
$M^*$ theorem's star condition and apply the cited common-region scaling
theorem.  Its exported hypothesis is affinity on each connected component.
For an open subset of $\mathbb R$ this is equivalent to local affinity:
components are open intervals, and affine formulas agree on overlapping
open neighborhoods, hence throughout each connected component.
\end{proofoutline}

This theorem derives the input sequence and limiting obstruction from native
correctness.  It lets a lower-bound argument analyze physical motion without
assuming a favorable trace class or conflating block time with execution time.
The toll is only $1/(2G)$ in total head-time, on the scale of half a scan of
a nominal $P/G$-cell storage footprint.  It is a bound on time outside affine
regions, not on distance travelled, and does not itself give a superlinear
time lower bound.

\begin{definition}[Bounded affine interval cover]
\label{def:motion-affine-cover}
A function $v:[0,H]\to\mathbb R$ has an affine interval cover of size $J$ if
there are $J$ endpoint pairs $\ell_j,r_j\in[0,H]$ such that their closed
intervals cover $[0,H]$ and $v$ is affine on each $[\ell_j,r_j]$.
Intervals may overlap or collapse; a pair with $\ell_j>r_j$ contributes an
empty interval.  No continuity or Lipschitz assumption is imposed on $v$.
\end{definition}
\formal{BoundedAffineCovers}{TransposeFollowup.AffineMotion.HasAffineIntervalCover}

\begin{lemma}[Pointwise closure]
\label{lem:motion-cover-closure}
If $v_i$ eventually have affine interval covers of a fixed finite size $J$
and $v_i(s)\to v(s)$ at every $s\in[0,H]$, then $v$ is locally affine on an
open subset of $[0,H]$ whose complement has measure zero.
\end{lemma}
\formal{BoundedAffineCovers}{TransposeFollowup.AffineMotion.local_affine_ae_of_eventual_bounded_covers}

\begin{proofoutline}
Pass to a subsequence on which all $2J$ endpoints converge.  Their limiting
closed intervals still cover $[0,H]$.  Away from the finitely many limiting
endpoints and the horizon endpoints, every point has a neighborhood lying
inside one interval for all sufficiently large indices.  On a smaller fixed
neighborhood the functions are affine.  Their coefficients converge because
the values at two distinct fixed points converge.  The pointwise limit is
affine there.
\end{proofoutline}

This closure lemma is the step that makes pointwise approximation sufficient
for the next native obstruction; no slope bound or uniform convergence of
the approximants is needed.

\begin{theorem}[Pointwise approximation obstruction]
\label{thm:motion-pointwise}
The native fast-correct hypotheses above are incompatible with the following
property.  For every strictly increasing $\phi:\mathbb N\to\mathbb N$ and
every input sequence $(w_i)$ eventually satisfying $w_i\in W_i$, there exist
a strictly increasing $\psi$, a finite $J$, and functions $v_{i,t}$ such that,
with $H=c+1/G$:
\begin{enumerate}
\item eventually every $v_{i,t}$ has an affine interval cover of size $J$ on
$[0,H]$;
\item for every head $t$ and every fixed $s\in[0,H]$,
\[
  \bigl|\widehat f_{\phi(\psi(i)),t}^{\,w_{\phi(\psi(i))}}(s)
                 -v_{i,t}(s)\bigr|\longrightarrow0.
\]
\end{enumerate}
Here $J$, the further subsequence, and the approximants may depend on the
chosen $\phi$ and the entire input sequence.  Only $J$ must be fixed along
that further subsequence.
\end{theorem}
\formal{PhysicalAffineApproximationObstruction}{TransposeFollowup.AffineMotion.no_subsequential_physical_pointwise_affine_fast_approximation}

\begin{proofoutline}
Choose the popular representatives and a convergent subsequence supplied by
Theorem~\ref{thm:motion-native-limits}, then apply the asserted approximation
property to precisely those choices.  The approximants converge pointwise to
the same physical limits.  Lemma~\ref{lem:motion-cover-closure} makes every
head locally affine outside a null set, contradicting
\eqref{eq:motion-physical-budget}.
\end{proofoutline}

Theorem~\ref{thm:motion-pointwise} reduces an exclusion proof to bounded-cover
approximation of every work head over the whole physical horizon.  Error
rates may vary with time, but convergence is required at \emph{every} fixed
time.  Regularity of a selected subroutine alone does not suffice.

\section{Concrete palettes and contracting recursion}
\label{sec:motion-applications}

We apply the motion obstruction to bounded words in explicit operation
palettes and to contracting trees of execution stages.  The palette proofs
supply the native implementations and their geometry.  The recursive result
requires the stated regularity condition on actual stages.

\subsection{The declared palettes}

All indices below start at zero.  Set $W=db$ and write the bank as
$x_{r,c,\ell}$, with $0\le r,c<d$ and $0\le\ell<b$, at flattened position
$rW+cb+\ell$.  A displayed formula specifies the output $y$ in terms of the
input $x$ to \emph{that operation}; a word applies its operations successively.
The following four symbol permutations are shared by the integer and ternary
palettes.  They preserve the order of the $b$ lanes inside each entry.
\begin{center}
\begin{tabular}{c l}
\hline
Operation & Exact output coordinate \\
\hline
$R$ & $y_{r,c,\ell}=x_{d-1-r,c,\ell}$ \\
$C$ & $y_{r,c,\ell}=x_{r,d-1-c,\ell}$ \\
$E$ & $y_{r,c,\ell}=x_{r,(c+r)\bmod d,\ell}$ \\
$e$ & $y_{r,c,\ell}=x_{r,(c-r)\bmod d,\ell}$ \\
\hline
\end{tabular}
\end{center}
Residues modulo $d$ are represented in $\{0,\ldots,d-1\}$.
The corresponding declarations are \texttt{rows}, \texttt{entries}, and the
two choices of \texttt{rotate}.
\formal{MatrixPermutationPalette}{TransposeFollowup.MatrixPermutationPalette.Operation,TransposeFollowup.MatrixPermutationPalette.value}
\formal{RawEntryRotation}{TransposeFollowup.RawEntryRotation.index_entry}
\formal{RawEntryReversal}{TransposeFollowup.RawEntryReversal.index_entry}

For the \emph{integer palette}, interpret each complete raw row as a
big-endian integer $X_r$ modulo $2^W$, and the complete bank as a big-endian
integer $X$ modulo $2^P$.  In addition to $R,C,E,e$, permit exactly the
following operations:
\begin{center}
\begin{tabular}{c l}
\hline
Operation & Output rule \\
\hline
Row sum & $Y_r=\sum_{a=0}^{r}X_a\pmod{2^W}$ \\
Row difference & $Y_0=X_0$, $Y_r=X_r-X_{r-1}\pmod{2^W}$ for $r>0$ \\
Row XOR prefix & $y_{r,c,\ell}=\bigoplus_{a=0}^{r}x_{a,c,\ell}$ \\
Row XOR difference & $y_{0,c,\ell}=x_{0,c,\ell}$;
 $y_{r,c,\ell}=x_{r,c,\ell}\oplus x_{r-1,c,\ell}$ for $r>0$ \\
Increment/decrement & $Y=X+1\pmod{2^P}$ or $Y=X-1\pmod{2^P}$ \\
\hline
\end{tabular}
\end{center}
Carries in row arithmetic can cross entry boundaries but not row boundaries.
A difference uses the preceding \emph{original} row; a prefix uses its
previous accumulated output.  Increment and decrement act on the whole bank.
\formal{IntegerMatrixPalette}{TransposeFollowup.IntegerMatrixPalette.Operation,TransposeFollowup.IntegerMatrixPalette.value}
\formal{RawRowArithmeticData}{TransposeFollowup.RawRowArithmetic.output_value}
\formal{RawRowXor}{TransposeFollowup.RawRowXor.bitValue_first,TransposeFollowup.RawRowXor.bitValue_next}
\formal{RawUnitArithmeticData}{TransposeFollowup.RawUnitArithmetic.output_value}

For the \emph{ternary palette}, the same coordinates hold values in
$\mathbb F_3$.  There are $b$ independent trit lanes per entry; these are not
integers represented in base three.  Along with $R,C,E,e$, the complete
additional palette is:
\begin{center}
\begin{tabular}{c l}
\hline
Operation & Exact output coordinate in $\mathbb F_3$ \\
\hline
$P$ & $y_{r,c,\ell}=\sum_{a=0}^{r}x_{a,c,\ell}$ \\
$p$ & $y_{r,c,\ell}=x_{r,c,\ell}-x_{r-1,c,\ell}$, with $x_{-1,c,\ell}=0$ \\
$Q$ & $y_{r,c,\ell}=\sum_{a=0}^{c}x_{r,a,\ell}$ \\
$q$ & $y_{r,c,\ell}=x_{r,c,\ell}-x_{r,c-1,\ell}$, with $x_{r,-1,\ell}=0$ \\
$D$ & $y_{r,c,\ell}=(-1)^{\operatorname{popcount}(r\mathbin{\&}c)}x_{r,c,\ell}$ \\
$S,s$ & $y_{a,\ell}=x_{(a+1)\bmod d^2,\ell}$ or
 $y_{a,\ell}=x_{(a-1)\bmod d^2,\ell}$ \\
\hline
\end{tabular}
\end{center}
In the last line $a=rd+c$ is the flattened \emph{entry} index.  In the phase
operation $\&$ is bitwise AND and $\operatorname{popcount}$ counts its set bits.
Binary inputs enter as trits $0,1$, and intermediate trits are retained
exactly.  Final emission applies $\pi(0)=0,\pi(1)=1,\pi(2)=0$; the theorem
therefore also excludes words whose final bank is already binary.
\formal{TernaryMatrixPalette}{TransposeFollowup.TernaryMatrixPalette.Operation,TransposeFollowup.TernaryMatrixPalette.value}
\formal{RawTernaryVertical}{TransposeFollowup.RawTernaryVertical.value_entry}
\formal{RawTernaryMatrixHorizontal}{TransposeFollowup.RawTernaryMatrixHorizontal.value_entry}
\formal{TernaryPhase}{TransposeFollowup.TernaryPhase.value_encoded_entry}
\formal{RawEntryShift}{TransposeFollowup.RawEntryShift.left_value_entry,TransposeFollowup.RawEntryShift.right_value_entry}
\formal{RawSymbolWord}{TransposeFollowup.RawSymbolWord.output}

The \emph{mixed binary palette} acts on $x_0,\ldots,x_{P-1}\in\mathbb F_2$.
A masked stride operation chooses a public mask $\mu:\mathbb N\to\{0,1\}$,
a stride $1\le S\le P$, a direction, and prefix or difference mode.  In the
forward direction it computes in increasing index order
\[
 y_j=x_j\oplus\bigl(\mu(j)z_j\bigr),\qquad
 z_j=\begin{cases}
 0,&j<S,\\
 y_{j-S},&j\ge S\text{ in prefix mode},\\
 x_{j-S},&j\ge S\text{ in difference mode}.
 \end{cases}
\]
Backward mode uses $j+S$ in place of $j-S$, computes in decreasing order,
and uses zero when $j+S\ge P$.  The remaining operations are full reversal
$y_j=x_{P-1-j}$ and rotation $y_j=x_{(j+R)\bmod P}$ with $0\le R\le P$.
These are all the operations in this palette.  The masks, strides, offsets,
and operation words may depend arbitrarily on size, but not on the input.
\formal{BinaryBankOperation}{TransposeFollowup.BinaryBankOperation.Operation,TransposeFollowup.BinaryBankOperation.Valid,TransposeFollowup.BinaryBankOperation.value}
\formal{MaskedStrideScan}{TransposeFollowup.MaskedStrideScan.buffer,TransposeFollowup.MaskedStrideScan.output}
\formal{BackwardMaskedStride}{TransposeFollowup.BackwardMaskedStride.value_recurrence}

\begin{theorem}[Bounded-word palette exclusions]
\label{thm:motion-palettes}
Let $d_i\to\infty$ and use the canonical $b_i,P_i$ above.  Fix any finite $K$.
For each of the three declared palettes, there is no sequence of words
$\omega_i$ of length at most $K$ eventually whose output equals raw transpose
eventually on sets $W_i$ of correctly encoded inputs satisfying
\eqref{eq:motion-entropy}.  For mixed binary words,
every stride and rotation must satisfy the stated bounds.  Public choices
may vary arbitrarily with $i$; there is no uniform-generation premise.
\end{theorem}
\formal{MatrixPaletteObstructions}{TransposeFollowup.IntegerMatrixPalette.no_high_entropy_transpose,TransposeFollowup.TernaryMatrixPalette.no_high_entropy_transpose}
\formal{RawMixedBinaryWordObstruction}{TransposeFollowup.RawMixedBinaryWord.no_high_entropy_transpose}

\begin{proofoutline}
Each word is compiled to one actual native execution.  The input is loaded
once, successive primitives run on the states produced by the preceding
primitives, and the output is emitted once.  Each primitive's contract
preserves the bank semantics, clears its scratch bank, and accounts for
rewinds and head restoration.

For the integer and ternary palettes, every primitive costs at most $24P$
and the input/output shell costs $7P$.  The executor has two work tapes, five
work symbols, $G=3$, and exactly $200(T+1)$ control states when its paid time
is $T$.  Thus $T\le(7+24K)P$ and its state-bit charge is $o(P)$.
Put $J_0=81^3\cdot9$.  At arbitrary prefix times and from arbitrary actual
stage states, every primitive has an affine-cover approximation with at most
$J_0$ pieces and unnormalized error at most $32(db+1)$.  The complete word,
including the shell and frozen tail, therefore has at most
$27J_0^K$ pieces and normalized error at most
\[
 \frac{32(db+1)K+4}{P}\longrightarrow0.
\]
The mixed binary executor has the same two tapes and five-symbol alphabet,
with $T\le(7+27K)P$ and $4(T+1)$ control states.  Strides need not be small
relative to $P$.  The geometry proof therefore selects subsequences according
to their normalized sizes: a vanishing stride has a vanishing-error
macroscopic scan approximation, while a stride bounded below as a fraction of
$P$ permits only boundedly many groups.  The resulting further subsequence
has bounded affine covers.  Masks change the data rule without changing the
paid head schedule.  The input/output phases are then added with their proved
error.  All three executors meet Theorem~\ref{thm:motion-pointwise}, which
contradicts the assumed correctness sets.
\end{proofoutline}
\formal{RawSymbolWord}{TransposeFollowup.RawSymbolWord.native_correct,TransposeFollowup.RawSymbolWord.resources,TransposeFollowup.RawSymbolWord.ticks_le}
\formal{RawSymbolWordApproximation}{TransposeFollowup.RawSymbolWord.scaled_approximation,TransposeFollowup.RawSymbolWord.state_bits_vanish}
\formal{IntegerMatrixGeometry}{TransposeFollowup.IntegerMatrixPalette.approximation_cover,TransposeFollowup.IntegerMatrixPalette.approximation_error}
\formal{TernaryMatrixGeometry}{TransposeFollowup.TernaryMatrixPalette.approximation_cover,TransposeFollowup.TernaryMatrixPalette.approximation_error}
\formal{RawMixedBinaryWord}{TransposeFollowup.RawMixedBinaryWord.fixed_resources}
\formal{RawMixedBinaryWordResources}{TransposeFollowup.RawMixedBinaryWord.bounded_ticks,TransposeFollowup.RawMixedBinaryWord.state_bits_vanish}
\formal{RawMixedBinaryWordRegularity}{TransposeFollowup.RawMixedBinaryWord.full_affine_approximation_subsequence}

Theorem~\ref{thm:motion-palettes} rules out these bounded-word construction
families without a caller-supplied geometry assumption.  A surviving word
construction must leave this class, for example by using unboundedly many
operations, input-dependent selection, or another primitive.  The asymptotic
exclusion is at $b=b(d)$; implementations of the primitives at other widths
do not by themselves extend it to those widths.

\subsection{Paid contracting stage clocks}

A contracting recursion may have unboundedly many stages.  To analyze it,
we assign every actual tick to a stage and bound the total duration of deep
stages.  A separate regularity condition will control each fixed stage's
physical motion.

\begin{definition}[Contracting stage clock]
\label{def:motion-clock}
Fix natural numbers $B,J$ and real parameters $0\le\rho<1$, $a\ge0$, and $C$.
For an execution of $T$ ticks at input size $P>0$, a clock consists of
nonnegative volumes $V(w)$ on addresses $w\in\{0,\ldots,B-1\}^{*}$, a label
$\ell:\{0,\ldots,T-1\}\to\{0,\ldots,B-1\}^{*}\times\{0,\ldots,J-1\}$,
and nonnegative integer endpoints $A(w,j)\le E(w,j)\le T$, such that
\begin{align}
 \sum_{b=0}^{B-1}V(wb)&\le\rho V(w), & V(\varnothing)&\le CP,\label{eq:motion-volumes}\\
 A(w,j)\le t<E(w,j)&\quad\Longleftrightarrow\quad\ell(t)=(w,j),
       &&(0\le t<T),\label{eq:motion-ownership}\\
 \sum_{j=0}^{J-1}\bigl(E(w,j)-A(w,j)\bigr)&\le aV(w).&&\label{eq:motion-own-budget}
\end{align}
The constants are fixed across the family, but the clock may depend on the
execution.  Empty stages are allowed.
\end{definition}
\formal{ContractingStageClock}{TransposeFollowup.RegularMotion.ContractingStageClock}

Equation~\eqref{eq:motion-ownership} partitions the global integer timeline.
Each own-stage is contiguous; the work of a node together with its descendants
need not be contiguous or nested.  The clock thus covers more schedules than
depth-first recursion.  Its volume $V(w)$ is an accounting quantity, not an
assumption of independently accessible memory.

Clock existence alone does not restrict linear-time executions when the
family constants may be chosen: if $T\le cP$, take $B=J=1$, $\rho=0$,
$a=1$, $C=c$, $V(\varnothing)=cP$, and one root stage $[0,T)$; give every
other node zero volume and an empty stage.  The substantive additional
hypothesis below is regularity of the stages marked good.

\begin{lemma}[Charged depth and exceptional work]
\label{lem:motion-clock-accounting}
For any such clock and $h\in\mathbb N$,
\begin{equation}
\label{eq:motion-depth}
 \#\{t<T:|\ell(t)_1|\ge h\}
 \le\frac{aV(\varnothing)\rho^h}{1-\rho},
 \qquad
 T\le\frac{aCP}{1-\rho}.
\end{equation}
Mark each stage good or exceptional, and let $E_{\rm bad}$ count all ticks
owned by exceptional stages.  The part of $[0,T/P)$ not covered by good
stages of depth less than $h$ has measure at most
\begin{equation}
\label{eq:motion-shallow}
 \frac{E_{\rm bad}}P+\frac{aC\rho^h}{1-\rho}.
\end{equation}
If exceptional ticks owned by each node $w$ number at most $\eta V(w)$,
where $\eta\ge0$, then
\begin{equation}
\label{eq:motion-node-charge}
 E_{\rm bad}/P\le\eta C/(1-\rho).
\end{equation}
\end{lemma}
\formal{StageClockAccounting}{TransposeFollowup.RegularMotion.stage_clock_deep_work_bound,TransposeFollowup.RegularMotion.stage_clock_total_time_bound}
\formal{ExceptionalStageClock}{TransposeFollowup.RegularMotion.regular_shallow_stage_omission,TransposeFollowup.RegularMotion.exceptionalTicks_node_budget}

\begin{proofoutline}
The total volume at depth $r$ is at most $\rho^rV(\varnothing)$.
Ownership identifies the number of ticks at one node with the sum of its
stage lengths.  Summing its own-work budget over depths at least $h$ gives
the geometric tail in \eqref{eq:motion-depth}.  The finite integer intervals
in \eqref{eq:motion-ownership}, divided by $P$, have measure exactly their
numbers of ticks divided by $P$.  Their disjointness is derived from the
ownership equivalence.  A tick outside the shallow good stages is either
exceptional or deep, giving \eqref{eq:motion-shallow}.  Charging exceptional
ticks instead by $\eta V(w)$ gives \eqref{eq:motion-node-charge}, with each
tick counted once.
\end{proofoutline}

The accounting lemma supplies the recursive obstruction's linear time bound
and controls the work lost when only finitely many depths are retained.  Its
per-node estimate will convert local exceptional-work budgets into a global
lower bound.

\subsection{Actual-stage regularity and the exceptional-time toll}

\begin{definition}[Regular sequences of stage occurrences]
\label{def:motion-stage-regularity}
Let $V_i,A_i,L_i$ be the volume, actual start tick, and duration of one fixed
stage across a sequence of executions, and let $h_i$ be one actual integer
head path.  These occurrences are regular if, for every strictly increasing
$\phi$ such that $L_{\phi(i)}>0$ for all $i$ and
$V_{\phi(i)}\to\infty$, there are a further strictly increasing $\psi$, a
finite $K$, and a function $v$ with an affine interval cover of size $K$ on
$[0,a]$ for which
\begin{equation}
\label{eq:motion-local-profile}
 \frac{h_{\phi(\psi(i))}\!\left(A_{\phi(\psi(i))}
       +\min\{\lfloor V_{\phi(\psi(i))}s\rfloor,L_{\phi(\psi(i))}\}\right)
       -h_{\phi(\psi(i))}(A_{\phi(\psi(i))})}
      {V_{\phi(\psi(i))}}
 \longrightarrow v(s)
\end{equation}
uniformly for $0\le s\le a$.  A sequence of flagged clocks is regular for a
head if this condition holds for every fixed address and stage, using its
true duration when the stage is good and duration zero otherwise.
\end{definition}
\formal{ActualStageRegularity}{TransposeFollowup.RegularMotion.HasRegularStageSubsequences}
\formal{ExceptionalRecursiveLimits}{TransposeFollowup.RegularMotion.RegularFlaggedClockHead}

The profile uses global ticks and freezes at the stage's actual end.  Its
cover size and limit may depend on the fixed stage, head, and subsequences;
there is no common piece bound across stages.  Missing, exceptional, and
bounded-volume occurrences impose no diverging positive-duration obligation.

Call a native family
\emph{flagged-recursive with allowance $(\eta,\varepsilon_i)$} if there exists
$N$ such that, for every strictly increasing $\phi$ with $\phi(i)\ge N$, every
sequence $w_i\in W_{\phi(i)}$, and every sequence of their actual halting
times $T_i$ with $P_{\phi(i)}>0$, there exist clocks and stage flags satisfying
\[
 E_{{\rm bad},i}/P_{\phi(i)}\le\eta+\varepsilon_{\phi(i)},
\]
and every work head is regular in Definition~\ref{def:motion-stage-regularity}.
The sequence $\varepsilon_i\to0$ is fixed for the family.  Clocks and flags
may depend on the execution sequence; no uniform algorithm for choosing them
is required.
\formal{NativeExceptionalRecursion}{TransposeFollowup.RegularMotion.HasFlaggedRecursiveExecutions}

\begin{theorem}[Exceptional-time toll for contracting recursion]
\label{thm:motion-exceptional-recursion}
Assume the native resources, state-bit bound, and high-entropy correct sets
of \eqref{eq:motion-control}--\eqref{eq:motion-entropy}.  Fix the clock
parameters $B,J,\rho,a,C$ as above.  Suppose every correct input eventually
has an actual halting transpose execution admitting a contracting stage
clock, and the family is flagged-recursive with allowance
$(\eta,\varepsilon_i)$, where $\varepsilon_i\to0$.  Then
\begin{equation}
\label{eq:motion-native-exception}
 \frac1{2G}\le k\eta,\qquad \eta\ge\frac1{2kG}.
\end{equation}
In particular, no such correct family can have every stage regular and zero
exceptional work.  The assertion holds for a fixed uniform machine as well.
If a linear-time bound is already assumed, the first inequality follows
without separately postulating clock existence for each correct execution.
\end{theorem}
\formal{NativeExceptionalRecursion}{TransposeFollowup.RegularMotion.native_flagged_recursion_budget,TransposeFollowup.RegularMotion.native_fast_flagged_recursion_budget,TransposeFollowup.RegularMotion.uniform_flagged_recursion_budget}
\formal{NativeRegularRecursion}{TransposeFollowup.RegularMotion.no_native_regular_recursion,TransposeFollowup.RegularMotion.no_uniform_regular_recursion}

\begin{proofoutline}
The clock gives $T_i\le aCP_i/(1-\rho)$.  Apply the popular-trace and native
limit theorem, discard a finite prefix, and invoke the family-level clock
hypothesis on the resulting representatives.

A diagonal subsequence supplies limits of the normalized stage endpoints,
volumes, and flags at every fixed node.  If a limiting good-stage interval
has positive length, its volume is bounded below by a positive multiple of
$P_i$, by the own-work budget.  Hence its volume tends to infinity and its
duration is eventually positive.  Definition~\ref{def:motion-stage-regularity}
then applies.  Transferring \eqref{eq:motion-local-profile} to the global
physical limit makes that limit locally affine on the interior of the stage
apart from finitely many breakpoints.  Zero-length limiting stages contribute
no time.

At each fixed depth there are finitely many stages.  Bound
\eqref{eq:motion-shallow} passes to their limiting intervals; sending the
depth to infinity leaves omitted measure at most $\eta$.  The countable union
of the locally affine stage regions therefore gives, for each head, an open
regular region missing at most $\eta$ physical time.  The frozen post-halting
tail is constant and is included in this region up to its endpoint.  Thus
\eqref{eq:motion-physical-budget} gives $1/(2G)\le k\eta$.
This also forces $k>0$, yielding the second inequality.  All-good flags and
$\eta=0$ prove the regular-recursion exclusion.
\end{proofoutline}
\formal{ExceptionalRecursiveLimits}{TransposeFollowup.RegularMotion.flagged_clock_limit_locally_affine}
\formal{ActualStageRegularity}{TransposeFollowup.RegularMotion.positive_stage_affine_of_regular_occurrences}
\formal{NativeRecursiveLimit}{TransposeFollowup.RegularMotion.native_flagged_limit_budget}

\begin{corollary}[Per-node exceptional budget]
\label{cor:motion-node-exception}
Under the same native correctness and clock assumptions, replace the global
exceptional allowance by the following family-level condition: in every
eligible execution sequence one can choose clocks and flags with the same
regularity and with at most $\eta V_i(w)$ exceptional ticks owned by each node
$w$, for a fixed $\eta\ge0$.  Then
\[
 \frac1{2G}\le k\frac{\eta C}{1-\rho},\qquad
 \frac{\eta C}{1-\rho}\ge\frac1{2kG}.
\]
\end{corollary}
\formal{NativeExceptionalRecursion}{TransposeFollowup.RegularMotion.HasNodeBudgetRecursiveExecutions,TransposeFollowup.RegularMotion.native_node_budget_recursion_bound}

\begin{proofoutline}
The exact node charge in \eqref{eq:motion-node-charge} produces a global
allowance $\eta C/(1-\rho)$ with error sequence identically zero.  Apply
Theorem~\ref{thm:motion-exceptional-recursion}.
\end{proofoutline}

Corollary~\ref{cor:motion-node-exception} gives a necessary lower bound on
per-node exceptional work.  Together with
Theorem~\ref{thm:motion-exceptional-recursion}, it excludes variants of the
geometric recursion in \cite[Section~13]{FirstPaper2026} whose actual stages
all have the regularity in Definition~\ref{def:motion-stage-regularity}.
It does not exclude the proposed non-affine M\"obius choreography: a
non-affine M\"obius map is affine on no nonempty open interval of its domain,
so a stage limit agreeing with such a map on an interval of positive length
has no finite affine cover.
The construction still has to realize its intended profiles on actual tapes.
The exceptional allowance must be at least $1/(2kG)$, a constant physical-time
toll corresponding asymptotically to $P/(2kG)$ ticks.  This toll does not grow
with the recursion depth.

These exclusions concern $O(P)$ transposition at canonical logarithmic
entry width.  Paid padding, proved next, shows that the first paper's
necessary $M^*$ condition can be supplied at additive linear cost to any
linear-time payload.  A construction must therefore both escape the excluded
motion classes and perform the required computation; a general lower bound
must identify an obligation that padding cannot satisfy independently.
Neither these restricted exclusions nor padding decides the unrestricted
transposition question.
